
\otazka{2. Bezpečnostní strategie}{Dokumenty České republiky a význam kyberbezpečnosti v nich. Bezpečnostní systém České republiky a role institucí v oblasti kyberbezpečnosti.}{BSSb1103}

\subsection*{2. Bezpečnostní strategie -- stručněji}
Tuhle otázku je dobré říct jako spojení dokumentů, institucí a praxe. Strategie řekne, co stát považuje za problém a jaký chce směr; bezpečnostní systém určí, kdo co dělá.

\textbf{1. Co je bezpečnostní strategie}
\begin{itemize}
  \item Strategie je politicko-koncepční dokument. Není to přímo zákon, ale nastavuje priority, cíle a způsob uvažování státu.
  \item Strategie vychází z analýzy bezpečnostního prostředí, hrozeb, rizik, národních zájmů a závazků v EU a NATO. Akční plán potom rozděluje úkoly, zákon ukládá povinnosti a metodika vysvětluje postup.
\end{itemize}

\textbf{2. Klíčové dokumenty ČR}
\begin{itemize}
  \item Bezpečnostní strategie ČR 2023 je nejširší rámec. Kyberbezpečnost chápe jako součást národní bezpečnosti, odolnosti, ochrany institucí, kritické infrastruktury a reakce na hybridní působení.
  \item Obranná strategie ČR 2023 řeší obranu státu, ozbrojené síly, odstrašení a vojenskou dimenzi. Kyber zde vystupuje hlavně jako doména obrany a multidoménových operací.
  \item Národní strategie kybernetické bezpečnosti 2026-2030 je specializovaný dokument pro kyberbezpečnost. Stojí na bezpečné strategické infrastruktuře, celospolečenské připravenosti a mezinárodní spolupráci.
\end{itemize}

\textbf{3. Bezpečnostní systém ČR}
\begin{itemize}
  \item Bezpečnostní systém ČR je síť orgánů, institucí, pravidel a procesů, které předcházejí hrozbám, reagují na ně a obnovují fungování státu.
  \item Patří sem prezident, Parlament, vláda, Bezpečnostní rada státu, ministerstva, kraje, obce, IZS, policie, zpravodajské služby, armáda a specializované úřady.
  \item V kyberbezpečnosti je hlavní civilní autoritou NÚKIB. Trestní rovinu řeší policie a státní zastupitelství, vojenskou kyberobranu Ministerstvo obrany, Vojenské zpravodajství a Armáda ČR.
\end{itemize}

\textbf{4. Praktický příklad}
\begin{itemize}
  \item Útok na nemocnici není jen IT problém. Nemocnice řeší obnovu, NÚKIB kybernetickou koordinaci, policie trestní rovinu, ÚOOÚ možný únik osobních údajů, zřizovatel kontinuitu péče a dodavatelé obnovu systémů.
  \item Strategický dokument říká, že odolnost zdravotnictví je důležitá; právní předpisy a metodiky pak určují, co mají subjekty prakticky dělat.
\end{itemize}

\subsection*{Kostra odpovědi}
\begin{enumerate}
  \item Vysvětlit funkci strategických dokumentů, akčních plánů, právních předpisů a metodik.
  \item Popsat klíčové dokumenty ČR: Bezpečnostní strategii 2023, Obrannou strategii 2023 a Národní strategii kybernetické bezpečnosti 2026-2030.
  \item Vymezit bezpečnostní systém ČR a jeho hlavní aktéry.
  \item Rozlišit instituce civilní kybernetické bezpečnosti, kyberkriminality, zpravodajství a kyberobrany.
  \item Ukázat na příkladu, jak se strategie promítne do právních povinností a praktické reakce.
\end{enumerate}

\subsection*{Strategie, akční plán, zákon a metodika}
Strategický neboli koncepční dokument určuje, jak stát chápe bezpečnostní problém, jakých dlouhodobých cílů chce dosáhnout a jakým směrem má postupovat. Bezpečnostní strategie je vrcholný strategický dokument, který popisuje bezpečnostní prostředí, zájmy státu, hlavní hrozby, priority a obecné způsoby reakce.

Názvy dokumentů se v praxi mohou překrývat, proto je důležitější jejich funkce než formální název. \textbf{Strategie} určuje problém, směr, cíle a priority. \textbf{Akční plán} převádí cíle do konkrétních úkolů, odpovědností a termínů. \textbf{Právní předpis} ukládá právně závazná práva, povinnosti a pravomoci. \textbf{Metodika} vysvětluje, jak mají být požadavky prakticky splněny.

Jednoduchý řetězec je: analýza prostředí a hrozeb - strategie - akční plán - právní, organizační a finanční opatření - vyhodnocení. Například strategie řekne, že je nutné zvýšit odolnost zdravotnictví proti ransomwaru; akční plán rozdělí úkoly; zákon a vyhlášky stanoví povinnosti; metodiky pomohou nemocnicím povinnosti splnit; vyhodnocení ukáže, zda se zlepšila detekce, reakce a obnova.

Strategické dokumenty mohou mít různý tematický záběr, časovou platnost, územní působnost a úroveň obecnosti. Bezpečnostní strategie ČR je široká a celostátní, Národní strategie kybernetické bezpečnosti se soustředí na jednu oblast, akční plán je konkrétnější a krajská nebo organizační strategie může mít užší územní či institucionální působnost.

\subsection*{Klíčové dokumenty ČR}
U zkoušky je lepší u každého dokumentu říct dvě věci: co je jeho funkce a jak přesně v něm kyberbezpečnost vystupuje. Není potřeba tvrdit, že všechny dokumenty řeší totéž. Bezpečnostní strategie je nejširší bezpečnostní rámec, Obranná strategie řeší obranu státu a vojenskou dimenzi, Národní strategie kybernetické bezpečnosti je specializovaný dokument pro kyberbezpečnost.

\textbf{Bezpečnostní strategie ČR 2023} je základní strategický dokument bezpečnostní politiky ČR. Kyberbezpečnost v ní vystupuje jako součást národní bezpečnosti, odolnosti státu a hybridního působení.
\begin{itemize}
  \item \textbf{Kyberprostor jako součást bezpečnostního prostředí:} strategie výslovně říká, že účinek hrozeb je násoben "překotným vývojem technologií, kybernetického a informačního prostoru i ekonomickými závislostmi".
  \item \textbf{Odolnost vůči nepřátelskému působení:} dokument uvádí, že "Česko musí být schopné odolávat nepřátelskému působení v kybernetické, informační, ekonomické a zpravodajské oblasti".
  \item \textbf{Konkrétní kyberhrozby:} v kyberprostoru protivníci útočí proti "kritické infrastruktuře a institucím", získávají citlivé informace a ohrožují strategické informační a komunikační systémy.
  \item \textbf{Kyberútok může mít strategickou váhu:} strategie počítá s tím, že "kybernetický útok velkého rozsahu může vést k aktivaci kolektivní obrany NATO".
  \item \textbf{Hybridní působení:} kybernetické, dezinformační, ekonomické a zpravodajské nástroje se podle strategie používají společně. Prakticky: kyberútok na instituci, únik dat, dezinformační kampaň a politický tlak mohou být části jedné operace.
  \item \textbf{Ochrana státu a společnosti:} zvyšování odolnosti se v dokumentu vztahuje i k ochraně kritické infrastruktury a k posilování informační a kybernetické bezpečnosti.
\end{itemize}

\textbf{Obranná strategie ČR 2023} řeší obranu státu, ozbrojené síly, odstrašení a spojenecké závazky. Kyber zde není obecná digitalizace státu, ale doména obrany, hybridního působení a multidoménových operací.
\begin{itemize}
  \item \textbf{Hybridní působení proti ČR:} strategie říká, že stát, společnost i ozbrojené síly jsou vystavovány "systematickému nepřátelskému hybridnímu působení".
  \item \textbf{Konkrétní hrozby:} jako hrozby pro ČR uvádí "kybernetické útoky, dezinformační kampaně, ekonomický nátlak a sabotážní, subverzní a zpravodajské aktivity".
  \item \textbf{Odolnost jako předpoklad obrany:} dokument zdůrazňuje, že základním nevojenským předpokladem obranyschopnosti je "posilování odolnosti státu, společnosti a občanů".
  \item \textbf{Kyberprostor jako vojenská doména:} strategie uvádí, že roste význam "vesmíru a kyberprostoru pro zajištění obranyschopnosti i pro fungování společnosti".
  \item \textbf{Multidoménové operace:} budoucí konflikty mají probíhat propojeně v pozemní, vzdušné, námořní, kybernetické a vesmírné doméně. Prakticky: kyberoperace může podporovat vojenské velení, rušit komunikaci protivníka nebo chránit vlastní systémy.
  \item \textbf{Schopnosti pro NATO a EU:} ČR má podle strategie budovat "schopnosti pro operace v kybernetickém prostoru" včetně kapacit využitelných ve prospěch NATO a EU.
  \item \textbf{Role NÚKIB v obranném systému:} strategie zmiňuje NÚKIB zejména u fungování kritické informační infrastruktury státu, ochrany utajovaných informací v informačních a komunikačních systémech, kryptografické ochrany a navigačních družicových systémů.
\end{itemize}

\textbf{Národní strategie kybernetické bezpečnosti 2026-2030} je hlavní specializovaný dokument pro kyberbezpečnost. Je dobré si zapamatovat její tři strategické oblasti a u každé jednu konkrétní formulaci.
\begin{itemize}
  \item \textbf{Bezpečná strategická infrastruktura:} strategie říká, že Česko bude posilovat zdroje a schopnosti "k proaktivní a účinné reakci na kybernetické hrozby", včetně budování kapacit kybernetické obrany.
  \item \textbf{Ochrana a odolnost:} dokument zdůrazňuje ochranu a odolnost strategické infrastruktury proti hrozbám z kyberprostoru i fyzického prostředí, včetně hybridního působení.
  \item \textbf{Kontinuita služeb:} když se útoku nepodaří zabránit, prioritou má být "zajištění kontinuity nebo včasné obnovy poskytování služeb". Prakticky: nemocnice, energetika nebo úřad musí mít zálohy, krizové postupy a plán obnovy, ne jen antivir.
  \item \textbf{Dodavatelské řetězce:} strategie výslovně řeší bezpečné a odolné dodávky, diverzifikaci a prověřování dodavatelů do strategické infrastruktury.
  \item \textbf{Celospolečenská připravenost a rozvoj:} dokument říká, že vysoká úroveň znalostí, schopností a kybernetické gramotnosti celé společnosti je zásadní předpoklad bezpečnosti kyberprostoru. Tím se kyberbezpečnost přesouvá i do vzdělávání, výzkumu, inovací a spolupráce veřejného a soukromého sektoru.
  \item \textbf{Mezinárodní spolupráce:} strategie uvádí, že bezpečnost Česka je úzce spjata s členstvím v NATO a EU "i v kyberprostoru".
  \item \textbf{Atribuce a reakce na škodlivé aktéry:} dokument počítá s asertivním vystupováním proti škodlivým aktérům, včetně atribuce útoků, diplomatické reakce a sankcí. Prakticky: veřejná atribuce APT31 proti MZV není jen technická zpráva, ale politicko-diplomatická reakce státu se spojenci.
\end{itemize}

Další dokumenty je vhodné zmínit jen stručně. \textbf{Strategie pro posílení odolnosti subjektů kritické infrastruktury 2026-2029} navazuje na all-hazards logiku CER a řeší odolnost základních služeb. Pro kyberobranu jsou relevantní aktuální dokumenty a koncepce resortu obrany; starší Strategie kybernetické obrany ČR 2018-2022 je historický dokument, ne aktuální strategie. Důležité jsou také Strategická koncepce NATO 2022 a unijní strategické dokumenty pro kybernetickou bezpečnost.

\subsection*{Bezpečnostní systém ČR}
Bezpečnostní systém ČR je soubor ústavních orgánů, státních institucí, ozbrojených sil, bezpečnostních sborů, záchranných složek, pravidel a procesů, kterými stát předchází hrozbám, reaguje na ně a obnovuje základní funkce. Jeho právní základ tvoří zejména Ústava, Listina základních práv a svobod, ústavní zákon o bezpečnosti České republiky, krizová legislativa a zvláštní bezpečnostní předpisy.

Hlavní aktéři:
\begin{itemize}
  \item \textbf{Prezident republiky} je hlava státu a vrchní velitel ozbrojených sil, ale není běžným řídicím orgánem civilní kyberbezpečnosti.
  \item \textbf{Parlament} přijímá zákony, kontroluje vládu a rozhoduje o některých zásadních bezpečnostních otázkách.
  \item \textbf{Vláda} je hlavní výkonný orgán bezpečnostní politiky.
  \item \textbf{Bezpečnostní rada státu} je stálý pracovní orgán vlády pro koordinaci bezpečnostních otázek a přípravu návrhů opatření.
  \item \textbf{Ústřední krizový štáb} je pracovní orgán pro koordinaci při závažných krizových situacích.
  \item \textbf{Ministerstva a jiné ústřední správní úřady} vykonávají své sektorové působnosti.
  \item \textbf{Kraje a obce} zajišťují územní krizové řízení a kontinuitu služeb.
  \item \textbf{Integrovaný záchranný systém} koordinuje záchranné a likvidační práce.
\end{itemize}

Kybernetická bezpečnost do systému vstupuje průřezově. Jeden incident může vyžadovat technickou reakci, právní posouzení, krizové řízení, komunikaci s veřejností, trestní řízení i mezinárodní koordinaci. Proto nestačí jeden úřad: NÚKIB je hlavní civilní autorita, ale ne jediný aktér.

\subsection*{Krizové stavy a stav kybernetického nebezpečí}
Krizové stavy nejsou totéž co každý závažný incident. Kybernetický útok může vést ke krizovému řízení, pokud má takový rozsah dopadů, že běžné nástroje nestačí. Vedle ústavních krizových stavů existuje i sektorový \textbf{stav kybernetického nebezpečí} podle zákona o kybernetické bezpečnosti. Stav kybernetického nebezpečí vyhlašuje Národní úřad pro kybernetickou a informační bezpečnost, tedy Úřad ve smyslu zákona č. 264/2025 Sb., na své úřední desce. Vyhlašuje se na nezbytně nutnou dobu, nejvýše na 30 dnů; prodloužení je možné jen v zákonných mezích tak, aby celková doba nepřekročila 60 dnů.

\subsection*{Role institucí v kyberbezpečnosti}
\textbf{NÚKIB} je hlavní civilní správní autorita pro kybernetickou bezpečnost. Vykonává regulaci, dohled, strategickou koordinaci a mezinárodní spolupráci, provozuje vládní CERT GovCERT.CZ, vydává výstrahy, varování a reaktivní protiopatření a koordinuje řešení významných incidentů ve své působnosti.

\textbf{CSIRT.CZ} je národní bezpečnostní tým provozovaný sdružením CZ.NIC. Pomáhá s koordinací incidentů zejména mimo vládní sektor, spolupracuje s poskytovateli, odbornou komunitou a zahraničními týmy. GovCERT.CZ a CSIRT.CZ se liší základním zaměřením, ale v praxi spolupracují.

\textbf{Policie ČR a státní zastupitelství} řeší trestní rovinu kybernetických incidentů: zajištění důkazů, identifikaci pachatele, trestní řízení a mezinárodní policejní i justiční spolupráci.

\textbf{Zpravodajské služby} sledují hrozby podle působnosti. BIS se zaměřuje na hrozby s dopadem na vnitřní bezpečnost ČR, ÚZSI na zahraniční informace a hrozby a Vojenské zpravodajství na obranu a vojenskou oblast. V kyberprostoru jsou důležité hlavně u státních aktérů, špionáže, strategické atribuce a varování.

\textbf{Kyberobrana} se musí odlišit od civilní kybernetické bezpečnosti. Civilní kybernetickou bezpečnost koordinuje především NÚKIB. Kybernetická obrana spadá do působnosti Ministerstva obrany, Vojenského zpravodajství a Armády ČR. Národní centrum kybernetických operací je spojeno s kybernetickou obranou a Vojenským zpravodajstvím. Velitelství kybernetických sil a informačních operací plánuje a řídí vojenské působení v kybernetickém a informačním prostoru. Není potřeba spekulovat o neveřejných schopnostech ani konkrétních ofenzivních operacích.

\textbf{Další aktéři} doplňují systém. Ministerstvo zahraničních věcí řeší kyberdiplomacii, mezinárodní reakce a veřejné atribuce. DIA je důležitá pro digitální služby státu, identitu a architekturu eGovernmentu. ČTÚ působí v elektronických komunikacích. NBÚ chrání utajované informace. Soukromý sektor provozuje velkou část infrastruktury a technologií. Akademická sféra zajišťuje vzdělávání, výzkum a odbornou podporu.

\subsection*{Praktické propojení strategie a institucí}
Příklad zdravotnictví ukazuje rozdíl mezi strategickým cílem, právní povinností, institucionální kompetencí a praktickou reakcí. Strategický dokument označí útoky na zdravotnictví za závažnou bezpečnostní hrozbu a stanoví cíl zvýšit odolnost nemocnic. Akční plán rozdělí úkoly mezi stát, zdravotnictví a regulátory. Zákon a vyhlášky stanoví povinnosti nemocnic, například bezpečnostní opatření, řízení rizik, hlášení incidentů a kontinuitu služby. Metodiky vysvětlí, jak se to má prakticky dělat.

Nemocnice pak zavede MFA, segmentaci, monitoring, školení, zálohy a plán kontinuity. Při incidentu zasahuje její technický tým a vedení, NÚKIB řeší kybernetickou rovinu a koordinaci ve své působnosti, policie trestní rovinu, ÚOOÚ případný únik osobních údajů, kraj nebo zřizovatel kontinuitu péče a dodavatelé obnovu systémů. Po incidentu se vyhodnotí, zda strategie a opatření fungovaly: jak rychle byl incident zjištěn, zda se nerozšířil, zda se obnovil provoz a zda nebyli nepřiměřeně ohroženi pacienti.

\priklad{APT31 a MZV ČR ukazuje strategickou a institucionální rovinu z jiné strany. Technická analýza a zpravodajské poznatky vedly k politickému rozhodnutí vlády o veřejné atribuci Čínské lidové republice. Do reakce vstoupily NÚKIB, zpravodajské služby, vláda, MZV a spojenci. Je to příklad, kdy kyberincident není jen technická událost, ale součást bezpečnostní politiky a kyberdiplomacie.}

\zkouska{Strategický dokument určuje, jak stát chápe bezpečnostní problém, jaké má dlouhodobé cíle a jakým směrem chce postupovat. Bezpečnostní strategie ČR 2023 je nejširší dokument bezpečnostní politiky, který vychází z analýzy prostředí, zdůrazňuje celovládní a celospolečenský přístup, odolnost a propojení vojenských, hybridních, kybernetických, ekonomických a informačních hrozeb. Obranná strategie ČR 2023 řeší obranu státu, ozbrojené síly, odstrašení a spojenecké závazky včetně vojenského působení v kybernetickém a informačním prostoru. Národní strategie kybernetické bezpečnosti 2026-2030 má tři hlavní oblasti: bezpečnou strategickou infrastrukturu, celospolečenskou připravenost a rozvoj a mezinárodní spolupráci a prosazování zájmů. Strategie není zákon: strategie určí směr, akční plán úkoly, právní předpis povinnosti a metodika praktický postup. Bezpečnostní systém ČR je soubor institucí, pravidel a procesů od prezidenta, Parlamentu, vlády a Bezpečnostní rady státu přes ministerstva, kraje, obce a IZS až po specializované úřady. V kyberbezpečnosti je hlavní civilní autoritou NÚKIB s GovCERT.CZ, vedle něj působí CSIRT.CZ, policie a státní zastupitelství, zpravodajské služby, Ministerstvo obrany, Vojenské zpravodajství, Armáda ČR, MZV, DIA, ČTÚ, NBÚ, soukromý sektor a akademická sféra. Klíčové je rozlišit civilní kybernetickou bezpečnost, trestní rovinu kyberkriminality a vojenskou kyberobranu.}

\zdroje{Přednáškový materiál \emph{2.) Bezpečnostní strategie (LONG)}; Bezpečnostní strategie ČR 2023; Obranná strategie ČR 2023; NÚKIB: Národní strategie kybernetické bezpečnosti 2026-2030; ústavní zákon o bezpečnosti České republiky; zákon o kybernetické bezpečnosti (č. 264/2025 Sb.).}

\newpage

